\documentclass[12pt,a4paper]{article}
\usepackage[T2A]{fontenc}
\usepackage[utf8]{inputenc}
\usepackage[russian]{babel}
\usepackage{caption}
\usepackage{amsmath, amssymb}
\usepackage{geometry}
\usepackage{tikz}
\usepackage{tikz-3dplot}
\usetikzlibrary{arrows.meta, calc, decorations.pathmorphing,
decorations.pathreplacing}
\geometry{margin=2cm}

\newcommand{\va}{\bar{a}}
\newcommand{\vb}{\bar{b}}
\newcommand{\vc}{\bar{c}}
\newcommand{\vi}{\bar{i}}
\newcommand{\vj}{\bar{j}}
\newcommand{\vk}{\bar{k}}
\newcommand{\pr}{\text{пр}}

\begin{document}

\begin{center}
  {\large\bfseries ЛЕКЦИЯ 8. СКАЛЯРНОЕ, ВЕКТОРНОЕ И СМЕШАННОЕ\\
  ПРОИЗВЕДЕНИЯ ВЕКТОРОВ}
\end{center}

\section*{8.1. Скалярное произведение векторов}

\begin{figure}[h]
  \centering
  \begin{minipage}{0.4\textwidth}
    \centering
    \begin{tikzpicture}[>=Stealth]
      \draw[->] (0,0) -- (3,0) node[below] {$A$};
      \node[below] at (1.5,0) {$\va$};
      \draw[->] (0,0) -- (1.8,2.2) node[right] {$B$};
      \node[left] at (0.7,1.3) {$\vb$};
      \draw (0.6,0) arc[start angle=0, end angle=51, radius=0.6];
      \node at (0.85,0.25) {$\alpha$};
      \node[below left] at (0,0) {$O$};
    \end{tikzpicture}
  \end{minipage}%
  \begin{minipage}{0.55\textwidth}
    Пусть $\va$ и $\vb$ --- ненулевые векторы.
    Приведём их к одному началу ---
    точка $O$. Пусть $\va = \overline{OA}$; $\vb = \overline{OB}$.
  \end{minipage}
  \caption*{Рисунок 8.1}
\end{figure}

\textbf{Определение.} Углом $\alpha$ между векторами $\va$ и $\vb$
называется наименьший угол $AOB$.

\textbf{Замечание.} $0 \leq \alpha \leq \pi$.

\medskip

\textbf{Определение.} \textit{Скалярным произведением векторов} $\va$
и $\vb$ называется число, равное:
$$
(\va,\, \vb) = |\va||\vb|\cos\alpha, \quad \text{где } \alpha =
\widehat{(\va,\vb)}.
$$

\textbf{Замечание.} $(\va, \vb) = |\va|\pr_{\va}\vb = |\vb|\pr_{\vb}\va$, т.\,к.
$$
\pr_{\vb}\va = |\va|\cos\alpha \;\Rightarrow\; \pr_{\vb}\va =
\frac{(\va,\vb)}{|\vb|}.
$$

\subsection*{Свойства скалярного произведения}

\begin{enumerate}
  \item $(\va, \vb) = (\vb, \va)$

  \item $(\va + \vb,\, \vc) = (\va, \vc) + (\vb, \vc)$

    \textit{Доказательство.}
    \begin{align*}
      (\va + \vb,\, \vc)
      &= |\vc|\pr_{\vc}(\va + \vb)
      = |\vc|\bigl(\pr_{\vc}\va + \pr_{\vc}\vb\bigr) \\
      &= |\vc|\pr_{\vc}\va + |\vc|\pr_{\vc}\vb
      = (\va, \vc) + (\vb, \vc). \qquad \text{ч.\,т.\,д.}
    \end{align*}

  \item $(\lambda\va,\, \vb) = \lambda(\va, \vb)$

  \item $(\va, \vb) = 0 \;\Leftrightarrow\; \va \perp \vb$, если $\va
    \neq \bar{0},\; \vb \neq \bar{0}$

  \item $(\va, \va) = |\va|^2 \;\Rightarrow\; |\va| = \sqrt{(\va, \va)}$
\end{enumerate}

\subsection*{Координатная запись скалярного произведения в данной
декартовой системе координат}

Пусть $\va = (x_1, y_1, z_1)$ и $\vb = (x_2, y_2, z_2)$.
\begin{align*}
  (\va, \vb)
  &= (x_1\vi + y_1\vj + z_1\vk,\; x_2\vi + y_2\vj + z_2\vk) \\
  &= (x_1\vi, x_2\vi) + (x_1\vi, y_2\vj) + (x_1\vi, z_2\vk)
  + (y_1\vj, x_2\vi) \\
  &\quad + (y_1\vj, y_2\vj) + (y_1\vj, z_2\vk)
  + (z_1\vk, x_2\vi) + (z_1\vk, y_2\vj) + (z_1\vk, z_2\vk) \\
  &= x_1x_2(\vi,\vi) + x_1y_2(\vi,\vj) + x_1z_2(\vi,\vk) \\
  &\quad + y_1x_2(\vj,\vi) + y_1y_2(\vj,\vj) + y_1z_2(\vj,\vk) \\
  &\quad + z_1x_2(\vk,\vi) + z_1y_2(\vk,\vj) + z_1z_2(\vk,\vk) \\
  &= x_1x_2 + y_1y_2 + z_1z_2.
\end{align*}

$$
\boxed{(\va, \vb) = x_1x_2 + y_1y_2 + z_1z_2.}
$$

\subsection*{Угол между векторами}

Пусть $\va = (x_1, y_1, z_1)$ и $\vb = (x_2, y_2, z_2)$.
$$
\cos\alpha = \frac{(\va,\vb)}{|\va||\vb|}, \quad \text{где } \alpha =
\widehat{(\va,\vb)}.
$$
$$
\cos\alpha =
\frac{x_1x_2 + y_1y_2 + z_1z_2}%
{\sqrt{x_1^2+y_1^2+z_1^2}\;\sqrt{x_2^2+y_2^2+z_2^2}}.
$$

\section*{8.2. Векторное произведение векторов}

\textbf{Определение.} Векторы называются \textit{компланарными}, если
они лежат в одной плоскости или в параллельных плоскостях.

\medskip

\textbf{Определение.} Тройку векторов, приведённых к одному началу,
называют \textit{упорядоченной}, если известно, какой из этих
векторов считается первым, какой вторым и какой третьим.

\medskip

\textbf{Определение.} Упорядоченная тройка некомпланарных векторов
называется \textit{правой}, если из конца третьего вектора кратчайший
поворот от 1-го вектора ко 2-му виден против часовой стрелки.

\begin{figure}[h]
  \centering
  \begin{minipage}{0.35\textwidth}
    \centering
    \begin{tikzpicture}[>=Stealth]
      \draw[->] (0,0) -- (2.5,0) node[below] {$\va$};
      \draw[->] (0,0) -- (1.6,1.4) node[right] {$\vb$};
      \draw[->] (0,0) -- (-0.3,2.8) node[above] {$\vc$};

      % внутренняя стрелка
      \draw[->, thick] (1.2,0.05)
      arc[start angle=-10, end angle=35, radius=1.1];

      % внешняя стрелка
      \draw[->, thick] (1.25,1)
      arc[start angle=55, end angle=-35, radius=0.65];

    \end{tikzpicture}
  \end{minipage}%
  \begin{minipage}{0.55\textwidth}
    $(\va,\vb,\vc)$ --- правая тройка\\[4pt]
    $(\vb,\va,\vc)$ --- левая тройка
  \end{minipage}
  \caption*{Рисунок 8.2}
\end{figure}

\textbf{Определение.} \textit{Векторным произведением векторов} $\va$
и $\vb$ называется вектор $\vc = [\va, \vb]$, удовлетворяющий
следующим условиям:
\begin{enumerate}
  \item $\vc \perp \va$, $\vc \perp \vb$;
  \item $|\vc| = |\va||\vb|\sin\alpha$, где $\alpha$ --- угол между
    векторами $\va$ и $\vb$;
  \item $(\va, \vb, \vc)$ --- правая тройка.
\end{enumerate}

\subsection*{Свойства векторного произведения}

\begin{enumerate}
  \item $[\va, \vb] = -[\vb, \va]$
  \item $[\va + \vb,\, \vc] = [\va, \vc] + [\vb, \vc]$
  \item $[\lambda\va,\, \vb] = \lambda[\va, \vb]$
  \item $[\va, \vb] = \bar{0} \;\Leftrightarrow\; \va \parallel \vb$,
    если $\va \neq \bar{0},\; \vb \neq \bar{0}$
  \item Если $\vc = [\va, \vb]$, то $|\vc| = S$, где $S$ --- площадь
    параллелограмма, построенного на векторах $\va$ и $\vb$.
\end{enumerate}

% Рисунок 8.3 — параллелограмм и вектор c
\begin{figure}[h]
  \centering
  \begin{tikzpicture}[>=Stealth]
    % parallelogram: base a, side b
    \coordinate (O) at (0,0);
    \coordinate (A) at (3,0);
    \coordinate (B) at (1,1.5);
    \coordinate (C) at (4,1.5);

    % fill hatching manually with lines
    \foreach \x in {0.15,0.4,...,2.9}{
      \draw[gray, thin] (\x,0) -- ++(1,1.5);
    }

    % white area around S
    \fill[white] (2.2,0.75) circle (0.25);

    % parallelogram outline
    \draw[->] (O) -- (A) node[below] {$\va$};
    \draw[->] (O) -- (B) node[left] {$\vb$};
    \draw (A) -- (C);
    \draw (B) -- (C);

    % area label
    \node at (2.2,0.75) {$S$};

    % vector c perpendicular, from origin upward
    \draw[->] (O) -- (0,3) node[left] {$\vc$};
  \end{tikzpicture}
  \caption*{Рисунок 8.3}
\end{figure}

\subsection*{Координатная запись векторного произведения в данной
декартовой системе координат}

Пусть $\va = (x_1, y_1, z_1)$ и $\vb = (x_2, y_2, z_2)$.
\begin{align*}
  [\va, \vb]
  &= [x_1\vi + y_1\vj + z_1\vk,\; x_2\vi + y_2\vj + z_2\vk] \\
  &= [x_1\vi, x_2\vi] + [x_1\vi, y_2\vj] + [x_1\vi, z_2\vk] + \ldots \\
  &= x_1x_2[\vi,\vi] + x_1y_2[\vi,\vj] + x_1z_2[\vi,\vk] + \ldots \\
  &= x_1y_2\vk - x_1z_2\vj + \ldots \\
  &= (y_1z_2 - z_1y_2)\vi - (x_1z_2 - x_2z_1)\vj + (x_1y_2 - y_1x_2)\vk \\
  &=
  \begin{vmatrix} \vi & \vj & \vk \\ x_1 & y_1 & z_1 \\ x_2 & y_2 & z_2
  \end{vmatrix}.
\end{align*}

$$
\boxed{[\va, \vb] =
  \begin{vmatrix} \vi & \vj & \vk \\ x_1 & y_1 & z_1 \\ x_2 & y_2 & z_2
\end{vmatrix}.}
$$

\section*{8.3. Смешанное произведение векторов}

\textbf{Определение.} \textit{Смешанным произведением} трёх векторов
$\va$, $\vb$ и $\vc$ называется число, равное
$$
(\va, \vb, \vc) = ([\va, \vb],\, \vc).
$$

\textbf{Теорема.} Пусть $V$ --- объём параллелепипеда, построенного
на векторах $\va$, $\vb$ и $\vc$. Тогда
$$
(\va, \vb, \vc) =
\begin{cases}
  +V, & \text{если } (\va, \vb, \vc) \text{ --- правая тройка;}\\
  -V, & \text{если } (\va, \vb, \vc) \text{ --- левая тройка;}\\
  0,  & \text{если векторы } \va,\, \vb \text{ и } \vc \text{ --- компланарны.}
\end{cases}
$$

% Рисунок 8.4 — параллелепипед
% TODO
\begin{figure}[h]
  \centering
  \begin{tikzpicture}[>=Stealth, scale=1.2]
    % Coordinates for the parallelepiped
    \coordinate (O) at (0,0);
    \coordinate (A) at (4,0);
    \coordinate (B) at (2.8,1.2);
    \coordinate (C) at (1.2,4);
    \coordinate (D) at (6.8,1.2);
    \coordinate (E) at (5.2,4);
    \coordinate (F) at (8,5.2);
    \coordinate (G) at (4,5.2);

    % Base hatching, less dense
    \begin{scope}
      \clip (O) -- (A) -- (D) -- (B) -- cycle;
      \foreach \x in {-1.5,-0.4,...,8} {
        \draw[thin, black!45] (\x,-0.5) -- ++(2.2,2.2);
      }
    \end{scope}

    % Dashed edges
    \draw[dashed, thick] (B) -- (D);
    \draw[dashed, thick] (B) -- (G);

    % Solid edges
    \draw[thick] (A) -- (D) -- (F) -- (G) -- (C) -- (E) -- (A);
    \draw[thick] (E) -- (F);

    % Vectors a, b, c
    \draw[->, thick] (O) -- (A) node[midway, below] {$\bar{a}$};
    \draw[->, thick] (O) -- (B) node[midway, above left=-2pt] {$\bar{b}$};
    \draw[->, thick] (O) -- (C) node[midway, right] {$\bar{c}$};

    % Vertical axis l
    \draw[->, thick] (0,-0.8) -- (0,6) node[left] {$\bar{l}$};

    % Horizontal line to indicate height level and right angle symbol
    \draw[thick] (0,4) -- (C);
    \draw (0,4.2) -- (0.2,4.2) -- (0.2,4);

    % Height h with brace on the other side
    \draw[
      decorate,
      decoration={brace, amplitude=8pt}
    ] (0,0) -- (0,4)
    node[midway, left=10pt] {$h$};

    % Angle phi between axis l and vector c
    \draw[thick] (0,1.15) arc[start angle=90, end angle=73, radius=1.15];
    \node at (0.32,1.65) {$\varphi$};

    % S_осн label and curved arrow pointing to the base
    \node (S) at (8.2,0.5) {$S_{\text{осн}}$};
    \draw[->, thick] (S.north)
    .. controls (7,3) and (5,3) ..
    (4.5,0.8);

  \end{tikzpicture}
  \caption*{Рисунок 8.4}
\end{figure}

\textbf{Доказательство.}
\begin{align*}
  (\va, \vb, \vc)
  &= ([\va, \vb],\, \vc)
  = |\vc|\,|[\va, \vb]|\cos\varphi \\
  &= |\vc|\cos\varphi \cdot S_{\text{осн}}
  = \pm h \cdot S_{\text{осн}} = \pm V.
\end{align*}

Если $(\va, \vb, \vc)$ --- правая тройка $\Rightarrow$ $\varphi$ ---
острый угол $\Rightarrow \cos\varphi > 0 \Rightarrow (\va, \vb, \vc) = +V$.

Если $(\va, \vb, \vc)$ --- левая тройка $\Rightarrow$ $\varphi$ ---
тупой угол $\Rightarrow \cos\varphi < 0 \Rightarrow (\va, \vb, \vc) = -V$.

Если векторы $\va$, $\vb$, $\vc$ --- компланарны $\Rightarrow
[\va,\vb] \perp \vc \Rightarrow ([\va,\vb], \vc) = 0 \Rightarrow
(\va, \vb, \vc) = 0$. \quad ч.\,т.\,д.

\subsection*{Координатная запись смешанного произведения в данной
декартовой системе координат}

Пусть $\va = (x_1, y_1, z_1)$, $\vb = (x_2, y_2, z_2)$ и $\vc = (x_3,
y_3, z_3)$.

$$
[\va, \vb] =
\begin{vmatrix} \vi & \vj & \vk \\ x_1 & y_1 & z_1 \\ x_2 & y_2 & z_2
\end{vmatrix}
=
\begin{vmatrix} y_1 & z_1 \\ y_2 & z_2
\end{vmatrix}\vi
-
\begin{vmatrix} x_1 & z_1 \\ x_2 & z_2
\end{vmatrix}\vj
+
\begin{vmatrix} x_1 & y_1 \\ x_2 & y_2
\end{vmatrix}\vk;
$$

\begin{align*}
  (\va, \vb, \vc) = ([\va,\vb],\, \vc)
  &=
  \begin{vmatrix} y_1 & z_1 \\ y_2 & z_2
  \end{vmatrix} x_3
  -
  \begin{vmatrix} x_1 & z_1 \\ x_2 & z_2
  \end{vmatrix} y_3
  +
  \begin{vmatrix} x_1 & y_1 \\ x_2 & y_2
  \end{vmatrix} z_3 \\
  &=
  \begin{vmatrix} x_3 & y_3 & z_3 \\ x_1 & y_1 & z_1 \\ x_2 & y_2 & z_2
  \end{vmatrix}
  =
  \begin{vmatrix} x_1 & y_1 & z_1 \\ x_2 & y_2 & z_2 \\ x_3 & y_3 & z_3
  \end{vmatrix}.
\end{align*}

$$
\boxed{(\va, \vb, \vc) =
  \begin{vmatrix} x_1 & y_1 & z_1 \\ x_2 & y_2 & z_2 \\ x_3 & y_3 & z_3
\end{vmatrix}.}
$$

Основное свойство смешанного произведения состоит в том, что
циклическая перестановка векторов не меняет его величины:
$$
(\va, \vb, \vc) = (\vb, \vc, \va) = (\vc, \va, \vb).
$$
В остальных случаях знак меняется на противоположный.

\section*{Примеры}

\subsection*{Пример 1}

Найти стороны и углы $\triangle ABC$, если $A(3;\,2;\,4)$;
$B(2;\,3;\,1)$; $C(4;\,-1;\,-2)$.

\textbf{Решение.}
\begin{align*}
  \overline{AB} &= (-1;\,1;\,-3), \quad AB = \sqrt{1+1+9} = \sqrt{11};\\
  \overline{AC} &= (1;\,-3;\,-6), \quad AC = \sqrt{1+9+36} = \sqrt{46};\\
  \overline{BA} &= (1;\,-1;\,3);\\
  \overline{BC} &= (2;\,-4;\,-3), \quad BC = \sqrt{4+16+9} = \sqrt{29}.
\end{align*}

$$
\cos\angle A
= \frac{(\overline{AB},\,\overline{AC})}{|\overline{AB}||\overline{AC}|}
= \frac{(-1)\cdot1 + 1\cdot(-3) + (-3)\cdot(-6)}{\sqrt{11}\cdot\sqrt{46}}
= \frac{-1-3+18}{\sqrt{506}}
= \frac{14}{\sqrt{506}};
$$

$$
\cos\angle B
= \frac{(\overline{BA},\,\overline{BC})}{|\overline{BA}||\overline{BC}|}
= \frac{1\cdot2 + (-1)\cdot(-4) + 3\cdot(-3)}{\sqrt{11}\cdot\sqrt{29}}
= \frac{2+4-9}{\sqrt{319}}
= -\frac{3}{\sqrt{319}}.
$$

$$
\angle A = \arccos\frac{14}{\sqrt{506}}; \qquad
\angle B = \arccos\!\left(-\frac{3}{\sqrt{319}}\right) = \pi -
\arccos\frac{3}{\sqrt{319}};
$$
$$
\angle C = \pi - \angle A - \angle B = \arccos\frac{3}{\sqrt{319}} -
\arccos\frac{14}{\sqrt{506}}.
$$

\subsection*{Пример 2}

Доказать, что диагонали четырёхугольника $ABCD$ перпендикулярны, если
$A(1;\,1;\,1)$; $B(3;\,2;\,-1)$; $C(2;\,3;\,-2)$; $D(-2;\,9;\,2)$.

\textbf{Решение.}
$$
\overline{AC} = (1;\,2;\,-3); \qquad \overline{BD} = (-5;\,7;\,3).
$$
$$
(\overline{AC},\,\overline{BD}) = -5 + 14 - 9 = 0 \;\Rightarrow\; AC \perp BD.
$$

\subsection*{Пример 3}

Найти площадь $\triangle ABC$, если $A(3;\,-1;\,2)$; $B(2;\,1;\,1)$;
$C(1;\,2;\,3)$.

\textbf{Решение.}
$$
\overline{AB} = (-1;\,2;\,-1); \qquad \overline{AC} = (-2;\,3;\,1).
$$
$$
\vc = [\overline{AB},\,\overline{AC}]
=
\begin{vmatrix} \vi & \vj & \vk \\ -1 & 2 & -1 \\ -2 & 3 & 1
\end{vmatrix}
=
\begin{vmatrix}2&-1\\3&1
\end{vmatrix}\vi
-
\begin{vmatrix}-1&-1\\-2&1
\end{vmatrix}\vj
+
\begin{vmatrix}-1&2\\-2&3
\end{vmatrix}\vk
= 5\vi + 3\vj + \vk;
$$
$$
|\vc| = \sqrt{25+9+1} = \sqrt{35} \;\Rightarrow\; S_{\triangle ABC} =
\frac{\sqrt{35}}{2}.
$$

\subsection*{Пример 4}

Найти объём пирамиды $ABCD$, если
$A(3;\,-1;\,0)$; $B(0;\,-3;\,3)$; $C(2;\,1;\,-1)$; $D(3;\,2;\,4)$.

\textbf{Решение.}
$$
\overline{AB} = (-3;\,-2;\,3); \qquad
\overline{AC} = (-1;\,2;\,-1); \qquad
\overline{AD} = (0;\,3;\,4).
$$
\begin{align*}
  (\overline{AB},\,\overline{AC},\,\overline{AD})
  &=
  \begin{vmatrix}
    -3 & -2 & 3 \\
    -1 &  2 & -1 \\
    0 &  3 & 4
  \end{vmatrix}
  =
  \begin{vmatrix}
    0 & -8 & 6 \\
    -1 &  2 & -1 \\
    0 &  3 & 4
  \end{vmatrix}
  =
  \begin{vmatrix}-8&6\\3&4
  \end{vmatrix}
  = -32 - 18 = -50.
\end{align*}
$$
V_{ABCD} = \frac{50}{6} = \frac{25}{3}.
$$

\subsection*{Пример 5}

Доказать, что точки $A(1;\,2;\,-1)$; $B(0;\,1;\,5)$; $C(-1;\,2;\,1)$;
$D(2;\,1;\,3)$ лежат в одной плоскости.

\textbf{Решение.}
$$
\overline{AB} = (-1;\,-1;\,6); \qquad
\overline{AC} = (-2;\,0;\,2); \qquad
\overline{AD} = (1;\,-1;\,4).
$$
$$
(\overline{AB},\,\overline{AC},\,\overline{AD})
=
\begin{vmatrix}
  -1 & -1 & 6 \\
  -2 &  0 & 2 \\
  1 & -1 & 4
\end{vmatrix}
=
\begin{vmatrix}
  5 & -1 & 6 \\
  0 &  0 & 2 \\
  5 & -1 & 4
\end{vmatrix}
= 0.
$$

Следовательно, векторы $\overline{AB}$, $\overline{AC}$ и
$\overline{AD}$ компланарны, а значит, точки $A$, $B$, $C$ и $D$
лежат в одной плоскости.

\end{document}
